\documentclass[11pt]{amsart}
\usepackage{amsmath,amssymb,amsthm}
\usepackage[margin=1.1in]{geometry}
\usepackage{array,longtable,booktabs}
\usepackage{hyperref}

\newtheorem{theorem}{Theorem}
\newtheorem{lemma}[theorem]{Lemma}
\newtheorem{proposition}[theorem]{Proposition}
\newtheorem{conjecture}[theorem]{Conjecture}
\theoremstyle{remark}
\newtheorem{remark}[theorem]{Remark}

\newcommand{\Inf}{\mathrm{Inf}}
\newcommand{\Var}{\mathrm{Var}}
\newcommand{\E}{\mathbb{E}}
\newcommand{\lab}{\mathrm{lab}}

\title{A filtration bound for the erasure entropies of Boolean functions}
\author{Yitzchak Shmalo}
\address{Einstein Institute of Mathematics, The Hebrew University of Jerusalem, Jerusalem, Israel}
\date{}

\begin{document}

\begin{abstract}
For a Boolean function $f$ and a coordinate $i$, the erasure entropy $E_i(f)$ is the conditional entropy of the event
$i\in U$ given $U\setminus\{i\}$, where $U$ is drawn from the spectral distribution $\hat f^2$. The Fourier
entropy-influence conjecture is equivalent to the existence of an absolute $c$ such that every non-constant $f$ has a
coordinate with $E_i(f)\le c\,\Inf_i(f)$. We bound $E_i$ by revealing the other coordinates one at a time and
conditioning only on the last revealed coordinate that the spectral sample contains. The resulting bound $L^\pi_i$
depends on the order $\pi$, lies between $E_i$ and the binary entropy of $\Inf_i$, and can be minimised over all orders
by a shortest-path recursion. We give a pivotality condition along the order under which $L^\pi_i=O(\Inf_i)$, an exact
formula for $E_i$ under block composition, and the bound $L^\pi_i/\Inf_i\le2+2^{-d}(2d+\log_2(1/\rho))$ for a
coordinate that enters only through one AND of $d+1$ variables and is pivotal with probability $\rho$. Conditioning on
the last $r$ revealed coordinates gives a decreasing family of bounds that reaches $E_i$ at $r=n-1$.
\end{abstract}

\subjclass[2020]{06E30, 42C10, 94A17}
\keywords{Fourier entropy-influence conjecture, Boolean functions, Fourier entropy, total influence}

\maketitle

\section{Introduction}\label{sec:intro}

For $f\colon\{-1,1\}^n\to\{-1,1\}$ write $f=\sum_S\hat f(S)\chi_S$, and let $U$ be a random subset with
$\Pr[U=S]=\hat f(S)^2$. The Fourier entropy is $H(f)=H(U)$ and the total influence is $I(f)=\E|U|=\sum_i\Inf_i(f)$, where
$\Inf_i(f)=\Pr[i\in U]$. The Fourier entropy-influence (FEI) conjecture of Friedgut and Kalai~\cite{FK96} asks for an
absolute $C$ with $H(f)\le C\,I(f)$ for every Boolean $f$. It is known for symmetric functions and read-once
decision trees~\cite{OWZ11} and for linear threshold functions~\cite{ShmLTF}, and it is preserved under suitable
compositions~\cite{OT13}. There are bounds in terms
of decision-tree and communication complexity measures~\cite{CKLS,WWW}; Arunachalam, Chakraborty, Kouck\'y, Saurabh
and de Wolf~\cite{ACKSW} bound $H(f)$ by twice the average unambiguous parity-certificate complexity. Kelman, Kindler,
Lifshitz, Minzer and Safra~\cite{KKLMS} prove a weak form for the min-entropy. For every Boolean $f$,
Han~\cite{Han} proved
\[
H(f)\le\Bigl(2+\frac{3}{\ln 2}\Bigr)I(f)+\sum_i\Inf_i(f)\log_2\frac{1}{\Inf_i(f)};
\]
Li and Han~\cite{LH25} lowered the constant in front of $I(f)$ to the sharp value $1/\ln 2$, and a form with $\Var(f)/\Inf_i(f)$ in place of $1/\Inf_i(f)$ in the logarithm is in~\cite{Shm4}.

One route to the conjecture goes through a single coordinate at a time. The drop of the entropy from $f$ to its two
restrictions at $x_i=\pm1$ splits into the erasure entropy (in the sense of Verd\'u and Weissman~\cite{VW08})
\[
E_i(f)=H\bigl(\mathbf 1[i\in U]\bigm| U\setminus\{i\}\bigr)
\]
where $\mathbf 1[\cdot]$ is the indicator, and a Jensen--Shannon term $J_i(f)\le(2/\ln 2)\Inf_i(f)$. The route goes back to Han~\cite[Remarks~3.1--3.2]{Han}; the
splitting is equation~(32) of Gonz\'alez, MacManus and Pereyra~\cite{GMP}. In
\cite{Shm1} it is shown that the conjecture holds if and only if there is an absolute $c$ such that every non-constant
Boolean function has a coordinate $i$ with
\begin{equation}\label{eq:good}
E_i(f)\le c\,\Inf_i(f),
\end{equation}
and in~\cite{Shm3} this quantity is studied directly. The difficulty is that $E_i$ conditions on the whole of
$U\setminus\{i\}$, an object with up to $2^{n-1}$ values. The trivial bound $E_i\le H(\mathbf 1[i\in U])$ loses a factor
$\log(1/\Inf_i)$, and~\eqref{eq:good} asks exactly to remove that factor at one coordinate.

Here we condition on less. Reveal the other coordinates in a fixed order and record only the position of the last
revealed coordinate that $U$ contains. This gives an upper bound $L^\pi_i$ for $E_i$ (Theorem~\ref{thm:witness}) with a
martingale form: at each step it compares how much of the average of $f$ over $x_i$ and how much of the derivative $D_if$ the
newly revealed coordinate explains. A two-sequence inequality turns a product condition along the order into the bound
$L^\pi_i=O(\Inf_i)$ (Theorem~\ref{thm:suff}), and the product condition has a direct reading in terms of resampling.
Section~\ref{sec:comp} gives an exact composition formula for $E_i$, Section~\ref{sec:private} the case of a private
AND, and Section~\ref{sec:windows} the windows of length $r$, which record the last $r$ revealed coordinates that $U$ contains. Section~\ref{sec:num} lists
exact values.

\section{The witness}\label{sec:witness}

Fix $i$ and write $g$ for the average of $f$ over $x_i$ and $h=D_if$ for its derivative, functions of the other coordinates
\[
g=\tfrac12\bigl(f|_{x_i=1}+f|_{x_i=-1}\bigr),\qquad h=\tfrac12\bigl(f|_{x_i=1}-f|_{x_i=-1}\bigr).
\]
They take values in $\{-1,0,1\}$, have disjoint supports, $\hat g(T)=\hat f(T)$ and $\hat h(T)=\hat f(T\cup\{i\})$ for
$T\not\ni i$, $\E h^2=\Inf_i(f)=:\delta$ and $\E g^2=1-\delta$. Put $p(T)=\hat f(T)^2$ and $q(T)=\hat f(T\cup\{i\})^2$ for
$T\subseteq[n]\setminus\{i\}$, and
\[
\Phi(a,b)=\eta(a)+\eta(b)-\eta(a+b)=(a+b)\,h_2\Bigl(\frac{b}{a+b}\Bigr),\qquad \eta(x)=-x\log_2x,
\]
where $h_2$ is the binary entropy. Then $E_i(f)=\sum_T\Phi(p(T),q(T))$, since $(U\setminus\{i\},\mathbf 1[i\in U])$ takes
the value $(T,0)$ with probability $p(T)$ and $(T,1)$ with probability $q(T)$. The function $\Phi$ is concave,
homogeneous of degree one and increasing in each argument; in particular it is superadditive.

Let $\pi=(j_1,\dots,j_m)$, $m=n-1$, be an order of the other coordinates, and let $S_k=\{j_1,\dots,j_k\}$. For
$T\subseteq[n]\setminus\{i\}$ let $\lab(T)=\max\{k:j_k\in T\}$, with $\lab(\emptyset)=0$, and put
\[
A_k=\sum_{\lab(T)=k}p(T),\qquad B_k=\sum_{\lab(T)=k}q(T),\qquad L^\pi_i(f)=\sum_{k=0}^m\Phi(A_k,B_k).
\]

\begin{theorem}\label{thm:witness}
For every Boolean $f$, coordinate $i$ and order $\pi$:
\begin{enumerate}
\item[(a)] $E_i(f)\le L^\pi_i(f)\le h_2(\Inf_i(f))$.
\item[(b)] With $\mathcal F_k$ the $\sigma$-algebra generated by $x_{j_1},\dots,x_{j_k}$,
\[
A_k=\E\bigl(\E[g\mid\mathcal F_k]-\E[g\mid\mathcal F_{k-1}]\bigr)^2,\qquad
B_k=\E\bigl(\E[h\mid\mathcal F_k]-\E[h\mid\mathcal F_{k-1}]\bigr)^2
\]
for $k\ge1$, and $A_0=(\E g)^2$, $B_0=(\E h)^2$.
\item[(c)] Let $a(R)=\sum_{T\subseteq R}p(T)$ and $b(R)=\sum_{T\subseteq R}q(T)$. Then $\min_\pi L^\pi_i=\lambda(R_*)$ for
the full set $R_*=[n]\setminus\{i\}$, where $\lambda(\emptyset)=\Phi(a(\emptyset),b(\emptyset))$ and
\[
\lambda(R)=\min_{j\in R}\Bigl(\lambda(R\setminus\{j\})+\Phi\bigl(a(R)-a(R\setminus\{j\}),\,b(R)-b(R\setminus\{j\})\bigr)\Bigr).
\]
\end{enumerate}
\end{theorem}

\begin{proof}
(a) The pair $(\lab(U\setminus\{i\}),\mathbf 1[i\in U])$ takes the value $(k,0)$ with probability $A_k$ and $(k,1)$ with
probability $B_k$, so $L^\pi_i=H(\mathbf 1[i\in U]\mid\lab(U\setminus\{i\}))$. Since $\lab(U\setminus\{i\})$ is a
function of $U\setminus\{i\}$, conditioning on it cannot give a smaller entropy than conditioning on
$U\setminus\{i\}$, and it cannot give a larger one than conditioning on nothing, which is $h_2(\delta)$.

(b) $\E[g\mid\mathcal F_k]=\sum_{T\subseteq S_k}\hat g(T)\chi_T$, so the $k$-th increment is
$\sum_{\lab(T)=k}\hat g(T)\chi_T$, whose squared norm is $A_k$ by Parseval. The same holds for $h$.

(c) The sets $\{T:\lab(T)\le k\}$ are the subsets of $S_k$, so $A_k=a(S_k)-a(S_{k-1})$ and $B_k=b(S_k)-b(S_{k-1})$. An
order is a maximal chain $\emptyset=S_0\subset S_1\subset\dots\subset S_m=R_*$, and $L^\pi_i$ is a sum of costs along
its edges. The recursion is the shortest path in this graph.
\end{proof}

The recursion costs $O(m2^m)$ operations per coordinate, against $m!$ orders. Part (b) says that $L^\pi_i$ follows the
Doob martingales of $g$ and of $h$ along the order and charges, at each step, the entropy of telling the two increments
apart. A coordinate is good along $\pi$ when $h$ is explained late, after most of $g$ has been explained.

\section{A sufficient condition}\label{sec:suff}

\begin{lemma}\label{lem:two}
Let $A_k,B_k\ge0$ for $k=0,\dots,m$, let $\delta=\sum_kB_k>0$, and suppose that for some $M\ge0$
\[
\Bigl(\sum_{k\ge K}A_k\Bigr)\Bigl(\sum_{k\le K}B_k\Bigr)\le M\delta^2\qquad\text{for every }K=0,\dots,m.
\]
Then $\sum_k\Phi(A_k,B_k)\le\bigl(1/\ln2+4+2\log_2(1+2M)\bigr)\delta$.
\end{lemma}

\begin{proof}
Since $a\log_2(1+b/a)\le b/\ln 2$, we have $\Phi(a,b)\le b/\ln2+\psi(a,b)$ with $\psi(a,b)=b\log_2(1+a/b)$. The function
$\psi$ is concave, homogeneous of degree one and increasing in each argument, so $\sum_{k\in G}\psi(A_k,B_k)\le
\psi(A_G,B_G)$ for every set $G$ of indices, where $A_G=\sum_{k\in G}A_k$. Put $\beta_k=\delta^{-1}\sum_{l\le k}B_l$,
which is nondecreasing, and $G_t=\{k:2^{-t-1}<\beta_k\le2^{-t}\}$ for $t\ge0$. Indices with $\beta_k=0$ have $B_k=0$ and
contribute nothing. For $k_t=\min G_t$ the hypothesis gives $A_{G_t}\le\sum_{k\ge k_t}A_k\le M\delta/\beta_{k_t}\le
2^{t+1}M\delta$, and $B_{G_t}\le\beta_{\max G_t}\delta\le2^{-t}\delta$. Hence
\[
\sum_{k\in G_t}\psi(A_k,B_k)\le2^{-t}\delta\log_2\bigl(1+2^{2t+1}M\bigr)\le2^{-t}\delta\bigl(2t+\log_2(1+2M)\bigr),
\]
and summing over $t$ with $\sum_t2^{-t}=2$ and $\sum_tt2^{-t}=2$ gives the claim.
\end{proof}

\begin{remark}\label{rem:two}
Under the hypothesis of Lemma~\ref{lem:two} one also has $\sum_k\Phi(A_k,B_k)\le\bigl(4/\ln2+\log_2(1+2M)\bigr)\delta$,
which is the smaller bound when $M\ge0.13$. Work in natural logarithms, put $\kappa=1+2M$ and let $\beta_k$ be as in
the proof of Lemma~\ref{lem:two}; indices with $B_k=0$ contribute only $A_k\ln1=0$. First,
$a\ln(1+b/a)\le b$. Second, for $B_k>0$ we have $\beta_k>0$, and $\ln x\le x-1$ at $x=\beta_k^2(A_k+B_k)/(\kappa B_k)$ gives
\[
B_k\ln\Bigl(1+\frac{A_k}{B_k}\Bigr)\le B_k\ln\kappa-2B_k\ln\beta_k+\frac{\beta_k^2A_k}\kappa,
\]
since $\beta_k^2\le1\le\kappa$. Third, $\sum_kB_k(-\ln\beta_k)=\delta\sum_k(\beta_k-\beta_{k-1})(-\ln\beta_k)\le
\delta\int_0^1(-\ln x)\,dx=\delta$, with $\beta_{-1}=0$. Fourth, with $\tau_l=\sum_{k\ge l}A_k\le M\delta/\beta_l$ (terms with $\beta_l=0$ vanish),
\[
\sum_kA_k\beta_k^2=\sum_l(\beta_l^2-\beta_{l-1}^2)\tau_l\le M\delta\sum_l\frac{(\beta_l-\beta_{l-1})(\beta_l+\beta_{l-1})}{\beta_l}
\le2M\delta.
\]
Adding, $\sum_k\Phi(A_k,B_k)\ln2\le\delta+\delta\ln\kappa+2\delta+2M\delta/\kappa\le(4+\ln\kappa)\delta$. In particular the constant
in Theorem~\ref{thm:suff} may be replaced by $4/\ln2+\log_2(1+2\max(M,1))$.
\end{remark}

For a set $R$ of coordinates other than $i$, let $x^R$ be $x$ with the coordinates in $R$ replaced by independent
uniform bits, and put
\begin{align*}
W(R)&=\E\,h(x)h(x^R)=\Pr[i\in U,\ U\cap R=\emptyset],\\
V(R)&=\E g^2-\E\,g(x)g(x^R)=\Pr[i\notin U,\ U\cap R\neq\emptyset].
\end{align*}
$W(R)$ is the correlation of the pivotality of $i$, with its sign, at $x$ and at $x^R$, and $V(R)$ is the influence of the block
$R$ on $g$. The two equalities are Parseval: $\E\,\varphi(x)\varphi(x^R)=\sum_{T\cap R=\emptyset}\hat\varphi(T)^2$ for every real function $\varphi$ of the coordinates other than $i$.

\begin{theorem}\label{thm:suff}
Let $\pi$ be an order and $R_k=R_*\setminus S_k$ the set of coordinates not yet revealed after step $k$. Then
$\sum_{l\le k}B_l=W(R_k)$ and $\sum_{l\ge k}A_l=V(R_{k-1})$. Consequently, if
\[
W(R_k)\,V(R_{k-1})\le M\,\Inf_i(f)^2\qquad\text{for }k=1,\dots,m,
\]
then $L^\pi_i(f)\le\bigl(1/\ln2+4+2\log_2(1+2\max(M,1))\bigr)\Inf_i(f)$.
\end{theorem}

\begin{proof}
$\lab(T)\le k$ means $T\subseteq S_k$, that is $T\cap R_k=\emptyset$, which gives the first identity; $\lab(T)\ge k$ means
$T\cap R_{k-1}\ne\emptyset$, which gives the second. For $K=0$ the product in Lemma~\ref{lem:two} is
$(1-\delta)\hat f(\{i\})^2\le\delta^2$, because $|\hat f(\{i\})|=|\E h|\le\Pr[h\ne0]=\delta$. Apply Lemma~\ref{lem:two}
with $\max(M,1)$.
\end{proof}

In words: $i$ is good along $\pi$ if at every stage either the value of $g$ is almost settled by the revealed
coordinates (the resampled block has small influence on $g$), or the revealed coordinates say little about whether $i$
is pivotal (resampling the rest destroys the pivotality of $i$). Because $W$ decreases and $V$ increases with $R$, a coarser
chain $R_*=R'_0\supset R'_1\supset\dots\supset\emptyset$ that the order refines suffices: if
$W(R'_{l+1})V(R'_l)\le M\Inf_i(f)^2$ for every $l$, the condition of Theorem~\ref{thm:suff} holds for every order that
reveals $R'_l\setminus R'_{l+1}$ in block $l$.

Write $\mathcal D=\{h\ne0\}$, $\mathcal L=\{g=-1\}$ and $\mathcal H=\{g=1\}$. Since $g\in\{-1,0,1\}$ and
$V(R)=\tfrac12\E(g(x)-g(x^R))^2$, every Boolean $f$ has $V(R)=\Pr[x\in\mathcal D,\,x^R\notin\mathcal D]+
4\Pr[x\in\mathcal L,\,x^R\in\mathcal H]$; if $f$ is monotone, then $h\ge0$ and $W(R)=\Pr[x\in\mathcal D,\,x^R\in\mathcal D]$.
For monotone $f$ the condition reads: the probability of staying in $\mathcal D$ times the probability of jumping
across it is $O(\Pr[\mathcal D]^2)$.

For example, let $f$ be the OR of $s$ disjoint ANDs of width $w$, let $i$ lie in the first AND, let $Q$ be the set of
the other $w-1$ variables of that AND, and let $P_0=(1-2^{-w})^{s-1}$ be the probability that no other AND is
satisfied. Take the chain $R_*\supset Q\supset\emptyset$, that is, reveal the other ANDs first and $Q$ last. Here
$x\in\mathcal D$ when no other AND is satisfied and every variable of $Q$ is true, so $\delta=\Inf_i(f)=P_02^{1-w}$ and
$W(Q)=P_04^{1-w}=\delta^2/P_0$; with $V(R_*)\le1$ the first product is at most $\delta^2/P_0$. A move from $\mathcal L$ to
$\mathcal H$ needs another AND to change, so resampling $Q$ cannot make it, and $V(Q)\le\Pr[x\in\mathcal D]=\delta$; with
$W(\emptyset)=\delta$ the second product is at most $\delta^2$. The condition therefore holds with $M=1/P_0$, which is
bounded in the usual regime $s\approx2^w\ln2$, where $P_0$ is close to $\tfrac12$.

\begin{conjecture}[filtration form]\label{conj:fwc}
There is an absolute $c$ such that every non-constant Boolean $f$ has a coordinate $i$ and an order $\pi$ of the other
coordinates with $L^\pi_i(f)\le c\,\Inf_i(f)$.
\end{conjecture}

By Theorem~\ref{thm:witness}(a) and~\cite{Shm1}, Conjecture~\ref{conj:fwc} implies the FEI conjecture with constant
$c+2/\ln2$: the class of Boolean functions is closed under restriction, and the one-coordinate step applies at the
coordinate given by the conjecture in every restriction. Conjecture~\ref{conj:fwc} is a strengthening of~\eqref{eq:good};
Section~\ref{sec:num} shows no counterexample among the functions we can evaluate exactly, and Section~\ref{sec:windows}
discusses where a single last element may not be enough.

\section{Composition}\label{sec:comp}

\begin{proposition}\label{prop:comp}
Let $X$ and $Z$ be disjoint sets of coordinates, let $\gamma\colon\{-1,1\}^X\to\{-1,1\}$ and
$F\colon\{-1,1\}\times\{-1,1\}^Z\to\{-1,1\}$ be Boolean, and put $f(x_X,z)=F(\gamma(x_X),z)$. Write $F(u,z)=F_0(z)+uF_1(z)$. For $i\in X$,
\[
\Inf_i(f)=\Inf_i(\gamma)\,\|F_1\|_2^2,
\]
\[
E_i(f)=\|F_1\|_2^2\Bigl(E_i(\gamma)-\Phi\bigl(\hat\gamma(\emptyset)^2,\hat\gamma(\{i\})^2\bigr)\Bigr)
+\sum_{T\subseteq Z}\Phi\Bigl(\bigl(\hat F_0(T)+\hat\gamma(\emptyset)\hat F_1(T)\bigr)^2,\ \hat\gamma(\{i\})^2\hat F_1(T)^2\Bigr).
\]
\end{proposition}

\begin{proof}
For $T_1\subseteq X$ and $T_2\subseteq Z$, $\hat f(T_1\cup T_2)=[T_1=\emptyset]\hat F_0(T_2)+\hat\gamma(T_1)\hat F_1(T_2)$.
For $i\notin T_1\ne\emptyset$ the pair $(p,q)$ at $T_1\cup T_2$ is $\hat F_1(T_2)^2$ times the pair of $\gamma$ at $T_1$, so by
homogeneity these terms sum to $\|F_1\|^2$ times the terms of $E_i(\gamma)$ with $T_1\ne\emptyset$. The terms with
$T_1=\emptyset$ form the last sum. The influence is the sum of $q$.
\end{proof}

The last sum depends on $\gamma$ only through $\hat\gamma(\emptyset)$ and $\hat\gamma(\{i\})$. When $\gamma$ is itself a composition this
can be iterated, which gives $E_i$ and $\Inf_i$ of iterated block compositions from the spectra of the gates; the
iteration for one gate applied to copies of another, carrying four numbers per leaf through each level, is in~\cite[Section~2]{Shm3}, where
it is used for towers of transitive gates. O'Donnell and Tan~\cite{OT13} prove a composition theorem for the
conjecture itself.

\section{A private AND}\label{sec:private}

\begin{proposition}\label{prop:private}
Let $i$ and a set $Y$ of $d\ge0$ other coordinates enter $f$ only through $\omega=\mathbf 1[x_i=x_j=1\text{ for all }j\in Y]$,
so that $f=f_0(v)+f_1(v)\,\omega$ with $v$ the remaining coordinates and $f_1\in\{0,\pm2\}$. Let $\delta=\Inf_i(f)=2^{-d}\Pr[f_1\ne0]>0$
and let $\pi$ reveal $v$ first and $Y$ last. Then
\[
2\bigl(1-2^{-d}\bigr)\le\frac{E_i(f)}{\delta}\le\frac{L^\pi_i(f)}{\delta}\le2\bigl(1-2^{-d}\bigr)+2^{-d}\log_2\frac{e\,2^d}{\delta}.
\]
\end{proposition}

\begin{proof}
In $\pm1$ notation $\omega=2^{-(d+1)}\sum_{S\subseteq Y\cup\{i\}}\chi_S$. Hence for $T_0\subseteq v$ and $S\subseteq Y$, the
coefficient of $f$ at $T_0\cup S\cup\{i\}$ is $2^{-(d+1)}\hat f_1(T_0)$, and at $T_0\cup S$ it is the same number when
$S\ne\emptyset$. Put $p_{T_0}=4^{-(d+1)}\hat f_1(T_0)^2$; then $\delta=\sum_Sq=2^d\sum_{T_0}p_{T_0}$. Every $T=T_0\cup S$ with $S\ne\emptyset$
has $p(T)=q(T)=p_{T_0}$ and contributes $\Phi(p_{T_0},p_{T_0})=2p_{T_0}$; there are $2^d-1$ such $S$ for each $T_0$, which gives
$2(1-2^{-d})\delta$, and the remaining terms of $E_i$ are nonnegative. In the witness, the label of $T_0\cup S$ with
$S\ne\emptyset$ is the position of the last element of $S$; the $l$-th element of $Y$ is the last element of $2^{l-1}$
sets $S$, so the atom at the position of that element has $A_k=B_k=2^{l-1}\sum_{T_0}p_{T_0}$ and contributes $2^l\sum_{T_0}p_{T_0}$, again $2(1-2^{-d})\delta$ in total. The
atoms of the sets $T=T_0$ carry $q$-mass $\nu=2^{-d}\delta$ and $p$-mass $\mu=1-2\delta+2^{-d}\delta$, so $\mu+\nu\le1$. By
superadditivity their contribution is at most $\Phi(\mu,\nu)\le\nu/\ln2+\nu\log_2((\mu+\nu)/\nu)\le\nu\log_2(e/\nu)$.
\end{proof}

Since $\delta=2^{-d}\Pr[f_1\ne0]$, the right-hand side is $2+2^{-d}\bigl(2d+\log_2e-2+\log_2(1/\Pr[f_1\ne0])\bigr)$, which
is less than $2+2^{-d}\bigl(2d+\log_2(1/\Pr[f_1\ne0])\bigr)$, so
the ratio tends to $2$ as the private AND grows, unless it is pivotal with probability below about $2^{-2^d}$. The same
computation applies to any private gadget whose Fourier coefficients at $S$ and $S\cup\{i\}$ agree in absolute value for
every non-empty $S$. Width-dependent bounds for every DNF, with no privacy assumption, are in~\cite{Shm3}.

\section{Windows}\label{sec:windows}

For $1\le r\le m$ let $\lab_r(T)$ be the set of the last $r$ positions of $\pi$ met by $T$ (all of them if $|T|<r$), and
\[
L^{\pi,r}_i(f)=H\bigl(\mathbf 1[i\in U]\bigm|\lab_r(U\setminus\{i\})\bigr).
\]
Then $L^{\pi,1}_i=L^\pi_i$, $L^{\pi,m}_i=E_i$, and
\[
E_i(f)\le L^{\pi,r}_i(f)\le L^{\pi,r-1}_i(f),
\]
because $\lab_{r-1}$ is a function of $\lab_r$, which is a function of $U\setminus\{i\}$. The masses of the atoms of $L^{\pi,r}_i$
are signed sums, by inclusion and exclusion, of resampling correlations of derivatives of order at most $r$: for
a set $\mathcal K$ of coordinates other than $i$ and a set $R$ of coordinates disjoint from it,
$\Pr[i\in U,\,\mathcal K\subseteq U,\,U\cap R=\emptyset]=\E\,D_{\mathcal K}h(x)\,D_{\mathcal K}h(x^R)$, where $D_{\mathcal K}$ is the product of the derivatives $D_j$ over $j\in\mathcal K$.

The ladder runs from Conjecture~\ref{conj:fwc} at $r=1$ to~\eqref{eq:good} at $r=m$. Window $1$ sees only the last element
of $U\setminus\{i\}$. When every variable lies in many overlapping certificates, the last element alone may carry
little information about whether $i\in U$, while a pair of elements from one certificate of $i$ may carry much more.
We do not know whether Conjecture~\ref{conj:fwc} holds, or whether it holds with some fixed window $r$ in place
of $1$.

\section{Exact values}\label{sec:num}

The statements of Theorems~\ref{thm:witness} and~\ref{thm:suff}, Lemma~\ref{lem:two} and Propositions~\ref{prop:comp},
\ref{prop:private} were checked with exact Walsh--Hadamard transforms. On $58$ random non-constant functions of $3$ to
$7$ variables (each value $-1$ with probability $\tfrac12$, $0.2$ or $0.05$), at every coordinate and for every order
when $n\le6$ and $150$ random orders when $n=7$: the bounds of Theorem~\ref{thm:witness}(a) hold to $3\cdot10^{-16}$,
the shortest path of Theorem~\ref{thm:witness}(c) agrees with the minimum over all orders to $3\cdot10^{-16}$, the
martingale form and the resampling identities hold exactly, the bound of Lemma~\ref{lem:two} holds with the smallest
admissible $M$, and $E_i\le L^{\pi,3}_i\le L^{\pi,2}_i\le L^{\pi,1}_i$ with $L^{\pi,m}_i=E_i$. The composition
formula on $60$ random compositions to $4\cdot10^{-16}$, and the bracket of Proposition~\ref{prop:private} on $80$ random
private ANDs. The table gives $\min_{i,\pi}L^\pi_i/\Inf_i$ and $\min_iE_i/\Inf_i$; $\mathrm{TRIBES}_{w\times s}$ is
the OR of $s$ disjoint ANDs of width $w$.
\[
\begin{array}{lcccccc}
 & \mathrm{AND}_6 & \mathrm{AND}_8 & \mathrm{AND}_{10} & \mathrm{TRIBES}_{2\times3} & \mathrm{TRIBES}_{2\times4} & \mathrm{TRIBES}_{2\times5}\\
\min_{i,\pi}L^\pi_i/\Inf_i & 2.2922 & 2.1048 & 2.0341 & 2.3160 & 3.0428 & 3.3892\\
\min_iE_i/\Inf_i & 2.2922 & 2.1048 & 2.0341 & 2.3160 & 3.0428 & 3.3892\\[4pt]
 & \mathrm{MAJ}_3 & \mathrm{MAJ}_5 & \mathrm{MAJ}_7 & \mathrm{MAJ}_9 & \mathrm{MAJ}_{11} & \mathrm{TRIBES}_{3\times3}\\
\min_{i,\pi}L^\pi_i/\Inf_i & 1.0000 & 1.5204 & 1.7848 & 1.9578 & 2.0844 & 2.5474\\
\min_iE_i/\Inf_i & 0 & 0 & 0 & 0 & 0 & 2.5474
\end{array}
\]
On AND and tribes the witness of the best order equals $E_i$ at the best coordinate. The functions $\mathrm{AND}_n$ are
the case $v=\emptyset$ of Proposition~\ref{prop:private}, and they sit just below its upper bound: at $n=10$,
$\min_{i,\pi}L^\pi_i/\Inf_i=2.034057$ against $2.034068$. On majority
$E_i=0$, because the spectrum lives on odd sets and $U\setminus\{i\}$ determines the parity of $|U|$, while the witness
pays for the last element; its growth along $n$ is slow and stays below $2.1$ at $n=11$.

The coordinate matters. For even $m$ let $\mathrm{IP}(y)=(-1)^{N(y)}$ on $\{-1,1\}^m$, where $N(y)$ is the number of $j\le m/2$
with $y_j=y_{j+m/2}=-1$ (the inner product function), and let $f(y,x_i)=\mathrm{IP}(y)$ for
$y\ne(1,\dots,1)$ and $f((1,\dots,1),x_i)=x_i$. For $m=2,4,6,8,10$ the coordinate $i$ has
$E_i/\Inf_i=2.672,\,5.224,\,7.387,\,9.429,\,11.439$ and $\min_\pi L^\pi_i/\Inf_i=2.950,\,5.351,\,7.426,\,9.439,\,11.442$,
while every coordinate $j$ of $y$ has $\min_\pi L^\pi_j/\Inf_j=0.891,\,1.699,\,1.929,\,1.983,\,1.996$. So $E_i/\Inf_i$
grows about linearly in $m$ at $i$ while every other coordinate stays below $2$, and the choice of the coordinate in
Conjecture~\ref{conj:fwc} cannot be dropped.

\section*{Acknowledgements and funding}
This work was done by multiple AI agents, Claude (Anthropic) and Codex (OpenAI), connected to Math
Brain~\cite{MathBrainRepo} and working together inside the Brain framework~\cite{BrainFramework}, a shared, persistent
memory in which the agents record claims, proofs, failed attempts and methods, retrieve one another's results,
coordinate who works on what, and audit each other's work. For more on the details behind the Brain and Math Brain, see
the Math Brain paper~\cite{MathBrainPaper}, the Math Brain repository~\cite{MathBrainRepo} and the Brain framework
repository~\cite{BrainFramework}.

The research presented in this paper was supported by the European Research Council (ERC) under the European Union's
Horizon Europe research and innovation programme (grant agreement No.~101041711), by the Simons Foundation as part of
the Collaboration on the Mathematical and Scientific Foundations of Deep Learning, by Heights Labs, by Convex Nexus
Capital, by the Israel Science Foundation (grant number 2258/19), and by the Israel Science Foundation (ISF Grant
4101/25). Computational resources were provided in part by Google Cloud through its education credits programme (credit
530550103).

\section*{Declaration of competing interest}
The author holds equity in Heights Labs and in Convex Nexus Capital and has a management role at Convex Nexus Capital.
These entities are included in the funding declaration above.

\section*{Declaration of generative AI and AI-assisted technologies}
The paper comes out of Math Brain, a project on the use of AI systems in research mathematics~\cite{BrainFramework,MathBrainRepo}. In November 2026 a
Claude Fable agent (Anthropic), connected to the Brain and working alongside Codex agents running OpenAI's GPT-6.1 on
the same problem, found the witness, its martingale form, the two-sequence condition and its resampling form, the
composition formula, the private-AND bound and the windows. Claude Opus 5.5 (Anthropic)
rederived each proof, wrote the numerical checks, prepared the manuscript and audited it. A further Claude agent
found the bound of Remark~\ref{rem:two} and checked it in Lean~4 with Mathlib (file in the code repository). The author directed the work and is responsible for the content of the paper.

\section*{Code and data availability}
Code and outputs: \url{https://github.com/yspennstate/fei-filtration-witness}, archived on Zenodo with doi:10.5281/zenodo.24125388.

The paper is part of the Math Brain project, whose public repository is
\url{https://github.com/yspennstate/math-brain-public}~\cite{MathBrainRepo}.

\appendix
\section{Notation}\label{app:notation}

Symbols are listed roughly in the order in which they first appear, and the last column gives that place. The letters $a$ and $b$ are the arguments of $\Phi$ and $\psi$ and also name the sums $a(R)$ and $b(R)$; the letter $l$ is an index in several proofs.

\begin{longtable}{@{}>{\raggedright\arraybackslash}p{0.19\textwidth}>{\raggedright\arraybackslash}p{0.54\textwidth}>{\raggedright\arraybackslash}p{0.2\textwidth}@{}}
\toprule
Symbol & Meaning & First appears \\
\midrule
\endhead
$f$, $n$ & a Boolean function $\{-1,1\}^n\to\{-1,1\}$, and its number of variables & Section~\ref{sec:intro} \\
$x$, $x_j$ & a point of $\{-1,1\}^n$, and its $j$-th coordinate & Section~\ref{sec:intro} \\
$i$, $j$ & coordinates; $i$ is the coordinate under study & Section~\ref{sec:intro} \\
$S$, $\hat f(S)$, $\chi_S$ & a set of coordinates, the Fourier coefficient of $f$ at $S$, and the character $\prod_{j\in S}x_j$ & Section~\ref{sec:intro} \\
$U$ & the spectral sample, a random set with $\Pr[U=S]=\hat f(S)^2$ & Section~\ref{sec:intro} \\
$H(\cdot)$, $H(f)$ & Shannon entropy in bits, and the Fourier entropy $H(U)$ & Section~\ref{sec:intro} \\
$\E$ & expectation; over a uniform point $x$ unless the variable is $U$ & Section~\ref{sec:intro} \\
$\Inf_i(f)$, $I(f)$ & the influence $\Pr[i\in U]$ of the coordinate $i$, and the total influence $\sum_i\Inf_i(f)$ & Section~\ref{sec:intro} \\
$C$ & the constant in the FEI conjecture $H(f)\le C\,I(f)$ & Section~\ref{sec:intro} \\
$\Var(f)$ & the variance of $f$ & Section~\ref{sec:intro} \\
$E_i(f)$ & the erasure entropy $H(\mathbf 1[i\in U]\mid U\setminus\{i\})$ & Section~\ref{sec:intro} \\
$\mathbf 1[\cdot]$ & the indicator of an event & Section~\ref{sec:intro} \\
$J_i(f)$ & the Jensen--Shannon term in the drop of the entropy from $f$ to its restrictions at $x_i=\pm1$ & Section~\ref{sec:intro} \\
$c$ & an absolute constant, in~\eqref{eq:good} and in Conjecture~\ref{conj:fwc} & Section~\ref{sec:intro} \\
$\pi$ & an order of the coordinates other than $i$ & Section~\ref{sec:intro} \\
$L^\pi_i(f)$ & the witness $\sum_k\Phi(A_k,B_k)=H(\mathbf 1[i\in U]\mid\lab(U\setminus\{i\}))$ & Section~\ref{sec:intro} \\
$D_j$ & the derivative in the direction $j$, $D_ju=\tfrac12(u|_{x_j=1}-u|_{x_j=-1})$ & Section~\ref{sec:intro} \\
$r$ & the length of a window, the number of last revealed positions that are recorded & Section~\ref{sec:intro} \\
$g$, $h$ & the average $\tfrac12(f|_{x_i=1}+f|_{x_i=-1})$ and the derivative $D_if$, functions of the coordinates other than $i$ & Section~\ref{sec:witness} \\
$\hat g(T)$, $\hat h(T)$ & the Fourier coefficients of $g$ and $h$ at $T$ & Section~\ref{sec:witness} \\
$\delta$ & $\Inf_i(f)=\E h^2$; in Lemma~\ref{lem:two}, $\sum_kB_k$ & Section~\ref{sec:witness} \\
$T$ & a subset of $[n]\setminus\{i\}$ & Section~\ref{sec:witness} \\
$p(T)$, $q(T)$ & $\hat f(T)^2$ and $\hat f(T\cup\{i\})^2$ & Section~\ref{sec:witness} \\
$\Phi(a,b)$ & $\eta(a)+\eta(b)-\eta(a+b)$, the entropy of telling a mass $a$ from a mass $b$ & Section~\ref{sec:witness} \\
$\eta(x)$, $h_2$ & $-x\log_2x$, and the binary entropy function & Section~\ref{sec:witness} \\
$j_k$, $m$ & the $k$-th coordinate of the order $\pi$, and $m=n-1$ & Section~\ref{sec:witness} \\
$k$ & a step of the order, $0\le k\le m$ & Section~\ref{sec:witness} \\
$S_k$ & $\{j_1,\dots,j_k\}$, the coordinates revealed after step $k$ & Section~\ref{sec:witness} \\
$\lab(T)$ & the label of $T$, the largest $k$ with $j_k\in T$, and $0$ for $T=\emptyset$ & Section~\ref{sec:witness} \\
$A_k$, $B_k$ & the $p$-mass and the $q$-mass of the sets with label $k$; in Lemma~\ref{lem:two}, any nonnegative numbers & Section~\ref{sec:witness} \\
$\mathcal F_k$ & the $\sigma$-algebra generated by $x_{j_1},\dots,x_{j_k}$ & Theorem~\ref{thm:witness} \\
$R$, $R_*$ & a set of coordinates other than $i$, and the full set $[n]\setminus\{i\}$ & Theorem~\ref{thm:witness} \\
$a(R)$, $b(R)$ & $\sum_{T\subseteq R}p(T)$ and $\sum_{T\subseteq R}q(T)$ & Theorem~\ref{thm:witness} \\
$\lambda(R)$ & the value of the shortest-path recursion at $R$; $\lambda(R_*)=\min_\pi L^\pi_i$ & Theorem~\ref{thm:witness} \\
$M$, $K$ & the constant in the product condition, and a cutoff $0\le K\le m$ & Lemma~\ref{lem:two} \\
$\psi(a,b)$ & $b\log_2(1+a/b)$ & proof of Lemma~\ref{lem:two} \\
$G$, $A_G$ & a set of indices, and $\sum_{k\in G}A_k$ & proof of Lemma~\ref{lem:two} \\
$\beta_k$ & $\delta^{-1}\sum_{l\le k}B_l$ & proof of Lemma~\ref{lem:two} \\
$t$, $G_t$, $k_t$ & a dyadic level, the indices $k$ with $2^{-t-1}<\beta_k\le2^{-t}$, and the least of them & proof of Lemma~\ref{lem:two} \\
$l$ & an index: of a sum, of a block of a coarser chain, or of an element of $Y$ & proof of Lemma~\ref{lem:two} \\
$\kappa$, $\tau_l$ & $1+2M$, and $\sum_{k\ge l}A_k$ & Remark~\ref{rem:two} \\
$x^R$ & $x$ with the coordinates in $R$ replaced by independent uniform bits & Section~\ref{sec:suff} \\
$W(R)$, $V(R)$ & $\E\,h(x)h(x^R)$, and $\E g^2-\E\,g(x)g(x^R)$, the influence of the block $R$ on $g$ & Section~\ref{sec:suff} \\
$\varphi$ & a real function of the coordinates other than $i$, in Parseval's identity & Section~\ref{sec:suff} \\
$R_k$ & $R_*\setminus S_k$, the coordinates not yet revealed after step $k$ & Theorem~\ref{thm:suff} \\
$R'_l$ & the sets of a coarser chain that the order refines & Section~\ref{sec:suff} \\
$\mathcal D$, $\mathcal L$, $\mathcal H$ & the sets $\{h\ne0\}$, $\{g=-1\}$ and $\{g=1\}$ & Section~\ref{sec:suff} \\
$s$, $w$ & the number and the width of the disjoint ANDs of a tribes function & Section~\ref{sec:suff} \\
$Q$, $P_0$ & the other $w-1$ variables of the AND that contains $i$, and the probability that no other AND is satisfied & Section~\ref{sec:suff} \\
$X$, $Z$, $z$ & the disjoint sets of coordinates of the inner and the outer function, and a point of $\{-1,1\}^Z$ & Proposition~\ref{prop:comp} \\
$\gamma$, $F$ & the inner Boolean function on $\{-1,1\}^X$, and the outer one on $\{-1,1\}\times\{-1,1\}^Z$ & Proposition~\ref{prop:comp} \\
$u$, $F_0$, $F_1$ & the first argument of $F$, and the parts in $F(u,z)=F_0(z)+uF_1(z)$ & Proposition~\ref{prop:comp} \\
$\hat\gamma$, $\hat F_0$, $\hat F_1$ & the Fourier coefficients of $\gamma$, $F_0$ and $F_1$ & Proposition~\ref{prop:comp} \\
$T_1$, $T_2$ & subsets of $X$ and of $Z$ & proof of Proposition~\ref{prop:comp} \\
$Y$, $d$ & the set of coordinates that enter only through the private AND, and its size & Proposition~\ref{prop:private} \\
$\omega$ & the private AND, the indicator that $x_i=1$ and $x_j=1$ for every $j\in Y$ & Proposition~\ref{prop:private} \\
$v$, $f_0$, $f_1$ & the remaining coordinates, and the functions with $f=f_0(v)+f_1(v)\,\omega$ & Proposition~\ref{prop:private} \\
$T_0$, $p_{T_0}$ & a subset of $v$, and $4^{-(d+1)}\hat f_1(T_0)^2$ & proof of Proposition~\ref{prop:private} \\
$\mu$, $\nu$ & the $p$-mass and the $q$-mass of the atoms of the sets $T=T_0$ & proof of Proposition~\ref{prop:private} \\
$e$ & the base of the natural logarithm & Proposition~\ref{prop:private} \\
$\lab_r(T)$, $L^{\pi,r}_i(f)$ & the last $r$ positions of $\pi$ met by $T$, and $H(\mathbf 1[i\in U]\mid\lab_r(U\setminus\{i\}))$ & Section~\ref{sec:windows} \\
$\mathcal K$, $D_{\mathcal K}$ & a set of coordinates other than $i$, and the product of the $D_j$ over $j\in\mathcal K$ & Section~\ref{sec:windows} \\
$\mathrm{AND}_n$, $\mathrm{MAJ}_n$, $\mathrm{TRIBES}_{w\times s}$ & the AND and the majority of $n$ variables, and the OR of $s$ disjoint ANDs of width $w$ & Section~\ref{sec:num} \\
$y$, $N(y)$, $\mathrm{IP}(y)$ & a point of $\{-1,1\}^m$, the number of $j\le m/2$ with $y_j=y_{j+m/2}=-1$, and the inner product function $(-1)^{N(y)}$ & Section~\ref{sec:num} \\
\bottomrule
\end{longtable}

\begin{thebibliography}{99}
\bibitem{FK96} E.~Friedgut, G.~Kalai, Every monotone graph property has a sharp threshold, Proc. Amer. Math. Soc.
124 (1996), 2993--3002.
\bibitem{ACKSW} S.~Arunachalam, S.~Chakraborty, M.~Kouck\'y, N.~Saurabh, R.~de~Wolf, Improved bounds on Fourier
entropy and min-entropy, ACM Trans. Comput. Theory 13 (2021), no.~4, 1--40, doi:10.1145/3470860; STACS 2020, arXiv:1809.09819.
\bibitem{CKLS} S.~Chakraborty, R.~Kulkarni, S.~V.~Lokam, N.~Saurabh, Upper bounds on Fourier entropy, Theoret.
Comput. Sci. 654 (2016), 92--112.
\bibitem{GMP} M.~J.~Gonz\'alez, P.~MacManus, M.~C.~Pereyra, Further evidence towards the Fourier
entropy-influence conjecture, arXiv:2606.00246, 2026.
\bibitem{Han} X.~Han, A new bound for the Fourier-entropy-influence conjecture, Combinatorica 45 (2025), no.~1,
Art.~4, doi:10.1007/s00493-024-00133-z; arXiv:2312.08271.
\bibitem{KKLMS} E.~Kelman, G.~Kindler, N.~Lifshitz, D.~Minzer, M.~Safra, Towards a proof of the
Fourier-entropy conjecture?, Geom. Funct. Anal. 30 (2020), 1097--1138.
\bibitem{LH25} P.~Li, G.~Han, Strengthening Han's Fourier entropy-influence inequality via an
information-theoretic proof, arXiv:2512.03117, 2025.
\bibitem{OT13} R.~O'Donnell, L.-Y.~Tan, A composition theorem for the Fourier entropy-influence conjecture, ICALP 2013,
Lecture Notes in Comput. Sci. 7965, 780--791, doi:10.1007/978-3-642-39206-1\_66.
\bibitem{OWZ11} R.~O'Donnell, J.~Wright, Y.~Zhou, The Fourier entropy-influence conjecture for certain classes of
Boolean functions, ICALP 2011, Lecture Notes in Comput. Sci. 6755, 330--341.
\bibitem{Shm1} Y.~Shmalo, The one-coordinate step in the Fourier entropy-influence conjecture,
\url{https://github.com/yspennstate/fourier-entropy-influence-one-coordinate}, doi:10.5281/zenodo.24125404, 2026.
\bibitem{Shm3} Y.~Shmalo, Good coordinates in the Fourier entropy-influence conjecture,
\url{https://github.com/yspennstate/fei-good-coordinates}, doi:10.5281/zenodo.24125398, 2026.
\bibitem{Shm4} Y.~Shmalo, A variance form of Han's bound on the Fourier entropy,
\url{https://github.com/yspennstate/fei-variance-han-bound}, doi:10.5281/zenodo.24125394, 2026.
\bibitem{ShmLTF} Y.~Shmalo, Fourier entropy of threshold functions and their compositions,
arXiv:2610.16069, 2026.
\bibitem{VW08} S.~Verd\'u, T.~Weissman, The information lost in erasures, IEEE Trans. Inform. Theory 54 (2008),
5030--5058.
\bibitem{WWW} A.~Wan, J.~Wright, C.~Wu, Decision trees, protocols, and the Fourier entropy-influence conjecture, ITCS 2014,
arXiv:1312.3003.
\bibitem{MathBrainRepo} Y.~Shmalo, \emph{Math Brain: public snapshot, code and tests}, 2026,
\url{https://github.com/yspennstate/math-brain-public}.
\bibitem{MathBrainPaper} Y.~Shmalo, \emph{Math Brain: shared, auditable memory for multi-agent mathematical research},
October 2026, \url{https://github.com/yspennstate/math-brain-public/tree/main/papers/math-brain-2026-10-24}, archived on Zenodo with doi:10.5281/zenodo.24150837.
\bibitem{BrainFramework} Y.~Shmalo, \emph{Brain Framework}, public software repository, 2026,
\url{https://github.com/yspennstate/brain-framework-public}.
\end{thebibliography}

\end{document}
